% Shared preamble for the literature/survey LaTeX documents.
% Used as the very first line of each document:  % Shared preamble for the literature/survey LaTeX documents.
% Used as the very first line of each document:  % Shared preamble for the literature/survey LaTeX documents.
% Used as the very first line of each document:  \input{preamble}
% (\input is a TeX primitive, so it is legal before \begin{document}.)
\documentclass[11pt]{article}
\usepackage[T1]{fontenc}
\usepackage{lmodern}
\usepackage[letterpaper,margin=1in]{geometry}
\usepackage{amsmath,amssymb}
\usepackage{array,booktabs}
\usepackage{enumitem}
\usepackage[most]{tcolorbox}
\usepackage{xcolor}
\usepackage[colorlinks=true,linkcolor=teal!60!black,urlcolor=teal!60!black,citecolor=teal!60!black]{hyperref}
\usepackage[backend=biber,style=numeric-comp,sorting=ynt,maxbibnames=8,maxcitenames=2,giveninits=true]{biblatex}
\addbibresource{references.bib}
\newtcolorbox{callout}{colback=black!3,colframe=black!25,boxrule=0.4pt,arc=2pt,
  left=6pt,right=6pt,top=4pt,bottom=4pt,breakable}
\setlist{topsep=3pt,itemsep=1pt,parsep=1pt,leftmargin=1.4em}
\setlength{\parindent}{0pt}\setlength{\parskip}{0.5em}
\sloppy

% (\input is a TeX primitive, so it is legal before \begin{document}.)
\documentclass[11pt]{article}
\usepackage[T1]{fontenc}
\usepackage{lmodern}
\usepackage[letterpaper,margin=1in]{geometry}
\usepackage{amsmath,amssymb}
\usepackage{array,booktabs}
\usepackage{enumitem}
\usepackage[most]{tcolorbox}
\usepackage{xcolor}
\usepackage[colorlinks=true,linkcolor=teal!60!black,urlcolor=teal!60!black,citecolor=teal!60!black]{hyperref}
\usepackage[backend=biber,style=numeric-comp,sorting=ynt,maxbibnames=8,maxcitenames=2,giveninits=true]{biblatex}
\addbibresource{references.bib}
\newtcolorbox{callout}{colback=black!3,colframe=black!25,boxrule=0.4pt,arc=2pt,
  left=6pt,right=6pt,top=4pt,bottom=4pt,breakable}
\setlist{topsep=3pt,itemsep=1pt,parsep=1pt,leftmargin=1.4em}
\setlength{\parindent}{0pt}\setlength{\parskip}{0.5em}
\sloppy

% (\input is a TeX primitive, so it is legal before \begin{document}.)
\documentclass[11pt]{article}
\usepackage[T1]{fontenc}
\usepackage{lmodern}
\usepackage[letterpaper,margin=1in]{geometry}
\usepackage{amsmath,amssymb}
\usepackage{array,booktabs}
\usepackage{enumitem}
\usepackage[most]{tcolorbox}
\usepackage{xcolor}
\usepackage[colorlinks=true,linkcolor=teal!60!black,urlcolor=teal!60!black,citecolor=teal!60!black]{hyperref}
\usepackage[backend=biber,style=numeric-comp,sorting=ynt,maxbibnames=8,maxcitenames=2,giveninits=true]{biblatex}
\addbibresource{references.bib}
\newtcolorbox{callout}{colback=black!3,colframe=black!25,boxrule=0.4pt,arc=2pt,
  left=6pt,right=6pt,top=4pt,bottom=4pt,breakable}
\setlist{topsep=3pt,itemsep=1pt,parsep=1pt,leftmargin=1.4em}
\setlength{\parindent}{0pt}\setlength{\parskip}{0.5em}
\sloppy

\begin{document}
\section*{ML Concepts --- a primer for the mathematician reading this literature}

\begin{callout}
Companion to the survey files in this directory. Audience: someone fluent in linear algebra, geometry, probability, and the basics of optimization, and comfortable with neural-network \emph{inference} (weights, activations, matrix multiplies) --- but \textbf{not} an ML specialist. The goal is to define, once and precisely, the ML vocabulary that the papers throw around without explanation. Where a term has a clean mathematical definition, it is given. Cross-references point to the companion \emph{Landscape} survey by part number.

\end{callout}

This is a living file: as later surveys hit new concepts, append them here rather than re-explaining inline.



\subsection*{1. The object of study: a decoder-only transformer}

Everything in this literature concerns \textbf{autoregressive decoder-only transformers} (GPT-style). Mathematically, such a model is a function


\[ \text{model}: \; (\text{sequence of tokens}) \longmapsto (\text{probability distribution over the next token}). \]


\begin{callout}
\textbf{The spaces in play} (named once, used throughout). Write $\mathcal V$ for the \textbf{vocabulary} (a finite set, $|\mathcal V| = n_{\text{vocab}}$). The model threads a handful of finite-dimensional real vector spaces:

\begin{itemize}
\item \textbf{token space} $T = \mathbb{R}^{\mathcal V}$ --- the free vector space on the vocabulary; a token is a basis vector (§1.1). Logits live in the dual $T^{*}$.
\item \textbf{feature space} $V = \mathbb{R}^{d_{\text{model}}}$ --- where each position's \emph{residual stream} vector lives (§1.3). The central space of the whole subject.
\item \textbf{position space} $P = \mathbb{R}^{n}$ --- one axis per token position in the context; only attention acts on this factor (§1.4).
\item \textbf{head space} $H = \mathbb{R}^{d_{\text{head}}}$ and \textbf{neuron space} $N = \mathbb{R}^{d_{\text{mlp}}}$ --- the (low- and high-dimensional) interiors of an attention head (§1.4) and an MLP (§1.5).
\end{itemize}

A theme worth stating at the outset: the model is built almost entirely from \textbf{linear maps between these spaces}. The handful of operations that instead refer to a \emph{chosen basis} of a space --- the softmax on $T^{*}$ (§1.2), the pointwise nonlinearity on $N$ (§1.5), the coordinate-wise normalization on $V$ (§1.6) --- are \emph{precisely} the ones that break basis-independence and create the "privileged bases" of §3. So we work with spaces and maps throughout, and descend to coordinates only for those few genuinely basis-dependent operations (where the coordinates are the point).

\end{callout}

\begin{callout}
\textbf{Adjoints vs. duals (a notational convention).} For a linear map $W:A\to B$ we keep two operations distinct, and never write a bare "$^\top$" for either:

\begin{itemize}
\item the \textbf{dual map} (a.k.a. transpose map) $W^{\vee}:B^{*}\to A^{*}$, $W^{\vee}(\phi) = \phi\circ W$ --- canonical, defined with \emph{no} inner product;
\item the \textbf{adjoint} $W^{*}:B\to A$ --- defined only once $A,B$ carry inner products, by $\langle Wa, b\rangle_B = \langle a, W^{*}b\rangle_A$ (the real-scalar specialization of the Hermitian adjoint).
\end{itemize}

The notation $W^{\vee}$ for the dual is \emph{not} universal --- texts also write ${}^{t}W$, $W'$, or even $W^{*}$; we use $W^{\vee}$ precisely so that $W^{*}$ can be reserved for the adjoint. The two coincide in coordinates with the standard inner product --- which is exactly why a transpose in the source papers is ambiguous. We resolve each occurrence: \textbf{$W^{*}$ always signals that a metric has been invoked; $W^{\vee}$ never does.}


Consequently a bare "$^{\top}$" survives in this document in \textbf{only one role}: the \emph{literal matrix transpose} inside an expression explicitly marked \emph{"coordinates"} or \emph{"matrix picture"} --- an optional, concrete realization of an object we have already named invariantly (e.g. the coordinate form $W_Q^{\top}W_K$ of the bilinear form $W_{QK}$, or $X A^{\top}$ for right-multiplying a position index). It never stands for a map between spaces. A quadratic form $x^{\top}M y$ likewise becomes the canonical pairing $\langle Mx, y\rangle$ with the metric $M:V\to V^{*}$ written as a map.

\end{callout}

\subsubsection*{1.1 Tokens, vocabulary, embedding, unembedding}
\begin{itemize}
\item \textbf{Token}: the atomic input symbol. Text is chopped by a \emph{tokenizer} into a sequence of tokens (roughly sub-words). The \textbf{vocabulary} $\mathcal V$ is the set of possible tokens, $|\mathcal V| = n_{\text{vocab}}$ (e.g. $\sim 50{,}000$); a token is a basis vector of \textbf{token space} $T = \mathbb{R}^{\mathcal V}$ (a \emph{one-hot} --- see box).
\item \textbf{Embedding}: a linear map $W_E : T \to V$ sending each token to its dense \textbf{embedding vector} in the feature space $V = \mathbb{R}^{d_{\text{model}}}$. (As a matrix it is $d_{\text{model}}\times n_{\text{vocab}}$.) Here $d_{\text{model}}$ --- a few hundred to many thousands --- is \emph{the} ambient dimension nearly every paper calls the residual-stream dimension.
\item \textbf{Unembedding}: a linear map $W_U : V \to T^{*}$ sending the final stream vector to a \textbf{logit covector}, which pairs with each token to score it.
\item Often $W_E, W_U$ are independent ("untied"); sometimes \textbf{tied} --- $W_U$ is the coordinate transpose of $W_E$, making the logit of token $y$ the inner product $\langle\operatorname{emb}(y),\,x\rangle_V$ of the stream with $y$'s embedding vector. That pairs two vectors \emph{of $V$}, so it needs an inner product on $V$: invariantly $W_U = W_E^{\vee}$ precomposed with the metric $V\cong V^{*}$. Even weight-tying thus silently picks a metric (cf. §1.3).
\end{itemize}

\begin{callout}
\textbf{On "one-hot".} "One-hot" is ML jargon for a \emph{standard basis vector}: under a fixed enumeration of the vocabulary, token $i$ is the vector $e_i \in \mathbb{R}^{n_{\text{vocab}}}$, all zeros but a single $1$ in coordinate $i$ ("exactly one coordinate is hot"; the name comes from digital logic). It is the canonical way to feed a \emph{categorical} variable --- a token drawn from an unordered finite set --- into linear algebra: encoding "token \#4217" as the \emph{scalar} $4217$ would impose a spurious order and metric, whereas the $e_i$ are linearly independent and mutually equidistant. It also makes embedding a literal table lookup, since $W_E e_i$ is simply the image of the basis vector $e_i$ (its "$i$-th column", in matrix terms); dually, the logit for token $i$ is the pairing $\langle \ell, e_i\rangle$ of the logit covector $\ell\in T^{*}$ with $e_i$. One nuance worth keeping: although a single token is \emph{one} $e_i$, the circuit papers freely take \emph{linear combinations} of one-hots (e.g. to reason about "a blend of two tokens"), so the real object is the token space $T = \mathbb{R}^{\mathcal V}$ with the tokens as a distinguished basis.

\end{callout}

\subsubsection*{1.2 Logits and softmax}
\begin{itemize}
\item \textbf{Logit}: the model's (unnormalized) score for a candidate next token. The \textbf{logit covector} $\ell = W_U\, x_{\text{final}} \in T^{*}$ assigns token $y$ the value $\ell_y = \langle \ell, e_y\rangle$ (its pairing with the one-hot $e_y$).
\item \textbf{Softmax}: the map $T^{*} \to \Delta(\mathcal V)$ onto probability distributions over the vocabulary,
\end{itemize}
\[ \operatorname{softmax}(\ell)_y = \frac{e^{\ell_y}}{\sum_{y'} e^{\ell_{y'}}}. \]   This is the first operation that refers to a \emph{basis} --- the token basis of $T^{*}$ --- but here the basis is genuinely meaningful (tokens are distinguished), unlike the arbitrary basis of $V$ (§1.3). Two invariances matter later: softmax is unchanged by \textbf{adding a constant} to every logit (so a logit is defined only modulo the all-ones covector $\mathbf 1$ --- an affine, not linear, object), and \textbf{temperature} scaling $\ell \mapsto \ell/\tau$ sharpens ($\tau\to 0$) or flattens ($\tau \to \infty$) the distribution.


\subsubsection*{1.3 The residual stream --- the central vector space}
The defining architectural choice, and the reason transformers are unusually amenable to linear-algebraic analysis. This subsection is worth reading slowly: most of the "circuits" program is downstream of the few exact statements below.


\textbf{Setup.} Fix a token position $i$. As the input flows through the network, that position carries a vector $x_i^{(k)} \in V$ (the feature space $\mathbb{R}^{d_{\text{model}}}$ of the overview) --- the \textbf{residual stream} at position $i$ after block $k$. The model is a sequence of $K$ \textbf{blocks} $B_1,\dots,B_K$ (each block is one attention layer or one MLP; in the finest-grained accounts each individual attention \emph{head} is its own block, since heads add independently --- §1.4). The stream is initialized to the embedding plus positional information (§1.7), \[ x_i^{(0)} \;=\; W_E\,(\text{token at position }i)\;+\;(\text{positional term}), \] and each block updates it through a \textbf{residual (skip) connection} --- it computes something and \emph{adds the result back} rather than replacing the stream: \[ x_i^{(k)} \;=\; x_i^{(k-1)} \;+\; B_k\!\big(x^{(k-1)}\big)_i. \] (Embedding and MLP blocks act position-wise --- $B_k(x^{(k-1)})_i$ depends only on $x_i^{(k-1)}$; attention is the \emph{only} block that reads other positions, which is why its argument is the whole $x^{(k-1)}$, not just position $i$.)


\textbf{Additivity, stated exactly.} Unrolling the recurrence telescopes into a sum: \[ x_i^{(K)} \;=\; x_i^{(0)} \;+\; \sum_{k=1}^{K} B_k\!\big(x^{(k-1)}\big)_i. \] The final residual vector is \emph{literally the sum} of the input embedding and every block's output: nothing is ever overwritten, only accumulated. The stream is a \textbf{running sum} and each block contributes one summand. This is the precise content of "additivity", and it is \textbf{exact and unconditional} --- it holds verbatim for every input, whatever the blocks compute.


\textbf{What is, and is not, linear.} A frequent misreading: "linear communication channel" does \emph{not} say the blocks are linear. Each $B_k$ is a genuinely \emph{nonlinear} function of the stream (softmax inside attention, ReLU/GeLU inside MLPs), so the step map $x^{(k-1)} \mapsto x^{(k)}$ is nonlinear. What is linear is (i) how the blocks' contributions \textbf{combine} --- by addition --- and (ii) how a contribution, once written, is \textbf{transported} to later blocks --- by the identity map, i.e. carried forward unchanged. The nonlinearity is confined \emph{inside} each block; the bus wiring the blocks together is linear. That separation is exactly what makes the circuits program possible.


\textbf{Logits are an additive sum of per-block contributions.} Since the unembedding $W_U$ is linear, applying it to the telescoped sum distributes over the terms (writing $\ell = W_U x_i^{(K)}$ for the logit vector of §1.2): \[ \ell \;=\; \underbrace{W_U x_i^{(0)}}_{\text{embedding's direct effect}} \;+\; \sum_{k=1}^{K} \underbrace{W_U\, B_k\!\big(x^{(k-1)}\big)_i}_{\text{block }k\text{'s direct effect on logits}}. \] Each summand is a vector in $\mathbb{R}^{n_{\text{vocab}}}$ that can be inspected on its own. Two standard techniques are immediate corollaries:

\begin{itemize}
\item \textbf{Direct logit attribution} --- to ask "which blocks pushed up the logit of output token $y$?", read off the $y$-component $\langle e_y,\, W_U B_k(\cdot)_i\rangle$ for each $k$.
\item \textbf{The logit lens} --- apply $W_U$ to an \emph{intermediate} $x_i^{(k)}$ to see the model's running guess at the next token. This is meaningful precisely because $x_i^{(k)}$ is a partial sum of the very same objects whose full sum produces the final logits.
\end{itemize}

\textbf{Reading and writing subspaces.} A block touches the stream through two linear maps: it \emph{reads} via an input map $W_{\text{in}}$ (e.g. $W_Q, W_K, W_V$ for attention, or the MLP's first map) and \emph{writes} via an output map $W_{\text{out}}$ (e.g. $W_O$, or the MLP's second map). What follows is what the stream's geometry says, exactly, about how blocks can and cannot interact --- and it is a good place to see which statements survive the gauge freedom and which are metric artifacts.


\emph{What a block produces and sees --- metric-free.} Locate the two maps by their relation to the stream space $V$: the write map $W_{\text{out}}$ has $V$ as its \textbf{destination}, the read map $W_{\text{in}}$ has $V$ as its \textbf{source}. Two statements then need no inner product:

\begin{itemize}
\item \textbf{Write.} The vector a block adds lies in $\operatorname{im}(W_{\text{out}}) \subseteq V$, its \textbf{write subspace}, of dimension $\le \operatorname{rank}(W_{\text{out}}) \le d_{\text{head}}$ (resp. $\le d_{\text{mlp}}$).
\item \textbf{Read.} The block depends on the stream vector $x$ only through $W_{\text{in}}x$, i.e. only through the class $\bar x \in V/\ker(W_{\text{in}})$. Inputs agreeing modulo $\ker(W_{\text{in}})$ are indistinguishable to it; equivalently, the block is \textbf{invariant to adding anything in $\ker(W_{\text{in}})$} to the stream.
\end{itemize}

Note the asymmetry in \emph{kind}, and that read and write live in \textbf{dual spaces}. The map-native \textbf{write} object is an honest \textbf{subspace of $V$} ($\operatorname{im}W_{\text{out}}$). The map-native \textbf{read} object is a \textbf{quotient of $V$} ($V \twoheadrightarrow V/\ker W_{\text{in}}$) --- equivalently, dually, the \textbf{subspace of covectors} $\operatorname{im}(W_{\text{in}}^{\vee}) \subseteq V^{*}$ that the block actually evaluates on the stream. So writes are vectors in $V$ and reads are covectors in $V^{*}$; a metric is exactly what one would need --- illegitimately --- to put them in the same space. Only the quotient / covector description is forced on us if we refuse a metric.


\emph{Turning the read side into a subspace of $V$ needs a metric --- and that is the catch.} To name the read side as a subspace of $V$ (rather than a quotient, or a subspace of $V^{*}$) one must choose a complement to $\ker(W_{\text{in}})$. The standard inner product supplies one, $V = \ker(W_{\text{in}}) \oplus \ker(W_{\text{in}})^{\perp}$, and then the block depends on $x$ only through the orthogonal projection $\Pi_{\text{read}}\,x$ onto $\ker(W_{\text{in}})^{\perp}$, since $W_{\text{in}}x = W_{\text{in}}(\Pi_{\text{read}}x)$. But §1.3 has \textbf{no privileged inner product}, and this complement is a \emph{metric artifact}. Under the gauge $x \mapsto Mx$ (so $W_{\text{in}}\mapsto W_{\text{in}}M^{-1}$) compare the three objects:

\begin{itemize}
\item the \textbf{kernel} transforms covariantly, $\ker \mapsto M\ker$ (a subspace of $V$);
\item its metric-free dual, the \textbf{annihilator} $(\ker W_{\text{in}})^{0}=\operatorname{im}(W_{\text{in}}^{\vee})\subseteq V^{*}$, transforms by the inverse \textbf{dual map} $(M^{\vee})^{-1}$;
\item the \textbf{orthogonal complement} $\ker^{\perp}$ transforms by the inverse \textbf{adjoint}, $(M\ker)^{\perp} = (M^{*})^{-1}(\ker)^{\perp}$ --- \emph{not} $M(\ker)^{\perp}$.
\end{itemize}

The complement and the annihilator agree only when $M$ is orthogonal ($M^{*}=M^{\vee}$ under the identification). So the kernel, the quotient $V/\ker$, and the annihilator are gauge-natural; the read subspace $\ker^{\perp}\subseteq V$ --- the "row space", obtained by dragging the annihilator across the metric isomorphism $V^{*}\cong V$ --- is not.


\emph{Non-interaction, formalized --- and gauge-invariant.} When does a writer $A$ leave a reader $B$ \textbf{completely} unaffected, on every input? $A$ contributes $W_{\text{out}}^A(\,\cdot\,)$ to the stream that $B$ reads through $W_{\text{in}}^B$, so $A$'s entire effect on $B$'s read is the composite $W_{\text{in}}^B W_{\text{out}}^A(\,\cdot\,)$. Hence \[ \underbrace{A \text{ never affects } B}_{\text{for every input}} \;\iff\; \operatorname{im}(W_{\text{out}}^A) \subseteq \ker(W_{\text{in}}^B) \;\iff\; \underbrace{W_{\text{in}}^B\, W_{\text{out}}^A = 0}_{\text{their virtual weight vanishes}}. \] All three are inner-product-free, and the virtual weight is \textbf{gauge-invariant} ($W_{\text{in}}^B M^{-1}\cdot M W_{\text{out}}^A = W_{\text{in}}^B W_{\text{out}}^A$). This is the \emph{correct} formalization of "$A$ and $B$ occupy orthogonal channels": not metric orthogonality of subspaces, but \textbf{image contained in kernel}. Degree of interaction is then the size of $W_{\text{in}}^B W_{\text{out}}^A$ --- its rank, or its norm once a metric is fixed.


\emph{Capacity, and why metric orthogonality returns as "interference".} The stream offers only $\dim V = d_{\text{model}}$ independent directions. Several writers' outputs can be linearly disentangled by a downstream reader exactly when their write subspaces are in \textbf{direct sum}, $\operatorname{im}W_{\text{out}}^{A_1}\oplus\cdots\oplus\operatorname{im}W_{\text{out}}^{A_k}\subseteq V$ (then projections $\Pi_j$ isolate each), which forces $\sum_j \dim\operatorname{im}W_{\text{out}}^{A_j} \le d_{\text{model}}$. But the model traffics in far more \emph{features} than $d_{\text{model}}$ (§3), so exact direct sums are impossible: write subspaces must overlap, and the virtual weights $W_{\text{in}}^B W_{\text{out}}^A$ between blocks that "should" be unrelated are forced \emph{nonzero}. That residual cross-talk \textbf{is} \emph{interference} (§3). Choosing the data-induced inner product (the one training commits to, §2), it is small precisely when the relevant write/read directions are \emph{nearly orthogonal} --- and the packing problem becomes "the feature Gram matrix is close to $I$" (§4, the \emph{Landscape} survey, Part III). So metric orthogonality is not fundamental, but it is the right \emph{approximate} language once a model has settled on a working inner product. (For attention, the overlap of different heads' write subspaces is itself a studied phenomenon, "attention superposition" --- \textcite{progress-on-attention}.)


\emph{A concrete channel.} The induction circuit above is one read/write channel made explicit. A \textbf{previous-token head} writes, into its write subspace $\operatorname{im}(W_O^{\text{prev}})$, a vector encoding "the token just before me was $t$". Layers later an \textbf{induction head}'s key map $W_K^{\text{ind}}$ is tuned to read that very subspace --- i.e. the virtual weight $W_K^{\text{ind}}W_O^{\text{prev}} \neq 0$ --- so it can match the current token against each earlier token's predecessor. The \emph{same} induction head's query map reads a different part of the stream (current-token identity), and the many other heads writing elsewhere are precisely those whose images $W_K^{\text{ind}}$ (approximately) annihilates: they share the bus without colliding with this channel.


So the residual stream is a shared memory bus of width $d_{\text{model}}$: blocks deposit into write subspaces and later blocks read out of the quotients dual to their read maps, with the virtual weight $W_{\text{in}}^B W_{\text{out}}^A$ the exact, gauge-invariant measure of how loudly $A$ speaks on $B$'s channel.


\textbf{Virtual weights.} Suppose an earlier block $A$ writes via $W_{\text{out}}^A$ and a later block $B$ reads via $W_{\text{in}}^B$. Because the identity carries $A$'s output forward untouched, the portion of $B$'s read that is \emph{due to} $A$ is governed by the single composite linear map \[ W_{\text{virtual}} \;:=\; W_{\text{in}}^B\, W_{\text{out}}^A, \] the \textbf{virtual weight} from $A$ to $B$. It is computable directly from the weights, with no reference to any input, and it exists \emph{even though $A$ and $B$ are not adjacent and the architecture contains no explicit weight connecting them} --- the residual connections supply the implicit identity wiring. Concretely: a \emph{previous-token head} writes a "the token before me was $t$" signal into some write subspace; several layers later an \emph{induction head} reads that subspace through its $W_K$; the virtual weight $W_K^{\text{ind}} W_O^{\text{prev}}$ is exactly the channel by which the first head speaks to the second, and analysing it is how \textcite{framework} reverse-engineers the induction circuit (§3, the \emph{Landscape} survey, Part I).


\textbf{Why this forces superposition.} The number of distinct signals the blocks collectively want to read and write vastly exceeds $d_{\text{model}}$ --- a single MLP already exposes $\sim\!4\,d_{\text{model}}$ neurons, each a potential feature. The stream is therefore a \textbf{bandwidth bottleneck}: many features must share the same $d_{\text{model}}$ dimensions, which is precisely the pressure toward \emph{superposition} (§3, and the \emph{Landscape} survey, Part I). Stream vectors forced to carry far more than $d_{\text{model}}$ signals are sometimes called "bottleneck activations" and are expected to be especially hard to interpret.


\textbf{No privileged basis.} Finally, because every read and write is an arbitrary linear map, one may apply \emph{any} automorphism $M \in GL(V)$ to the entire stream and absorb it into the weights --- replace each read map $W_{\text{in}}$ by $W_{\text{in}} M^{-1}$ and each write map $W_{\text{out}}$ by $M\,W_{\text{out}}$ --- for a bit-identical function. The stream thus has \textbf{no privileged basis}: its individual coordinates carry no intrinsic meaning, and the space comes with a $GL(V)$ \textbf{gauge symmetry}. (A pointwise nonlinearity, as on MLP neurons, \emph{would} break this and single out a basis --- §1.5.) The consequences of this gauge freedom --- in particular that within-stream quantities like cosine similarity are not canonically defined --- are the subject of the \emph{Landscape} survey, Part III.


\subsubsection*{1.4 Attention}
The block that moves information \emph{between} token positions, and the \textbf{only} one acting on the position factor $P$ (embedding and MLP act on the feature space $V$ alone, position by position). Assembling all positions, the residual stream at a layer is an element $X \in V \otimes P$ --- concretely the $d_{\text{model}}\times n$ matrix whose $j$-th column is the stream vector $x_j$ at position $j$.


\textbf{One head} is four linear maps through the low-dimensional head space $H = \mathbb{R}^{d_{\text{head}}}$ (with $d_{\text{head}} \ll d_{\text{model}}$): \[ W_Q, W_K, W_V : V \to H, \qquad W_O : H \to V. \] The head computes, in three steps:


\begin{enumerate}
\item \textbf{Scores (a bilinear form on $V$, landing in the dual).} Query and key maps send the stream to $H$-valued fields $W_Q X, W_K X$, and the score of destination $i$ attending to source $j$ is the \textbf{head-space inner product} $\langle W_Q x_i,\, W_K x_j\rangle_H$ (an \emph{adjoint} pairing on $H$). Invariantly this is a single learned \textbf{bilinear form} on $V$, \[ W_{QK} : V \to V^{*}, \qquad W_{QK} = W_Q^{\vee}\, g_H\, W_K \quad(\text{coordinates: } W_Q^{\top}W_K), \] where $g_H:H\xrightarrow{\sim}H^{*}$ is the inner product on $H$ and $W_Q^{\vee}$ the dual map. The score is then the \textbf{canonical pairing} $\langle W_{QK}\,x_j,\; x_i\rangle$ of a covector with a vector --- needing \textbf{no inner product on $V$}. The lone transpose (in the coordinate form) marks precisely that $W_{QK}$ lands in the dual $V^{*}$; and the one metric used, $g_H$ on the \emph{head} space, is \textbf{internal} --- a change of basis on $H$ rescales $W_Q, W_K$ oppositely and leaves $W_{QK}$ fixed. The form need not be symmetric (query- and key-roles differ), so it distinguishes "$i$ attends to $j$" from "$j$ attends to $i$". Evaluated on every pair of positions it gives the score array $S$ --- itself a bilinear form on the position space $P$.
\end{enumerate}

\begin{enumerate}
\item \textbf{Pattern (the only nonlinearity).} Normalize each row of $S$ (scaled by $1/\sqrt{d_{\text{head}}}$, and with a \emph{causal mask} forbidding attention to future positions, $j>i$) through a softmax: \[ A \;=\; \operatorname{rowsoftmax}\!\big(S/\sqrt{d_{\text{head}}}\;+\;\text{mask}\big) \;\in\; \operatorname{End}(P). \] $A$ is \textbf{row-stochastic} (rows are probability vectors) and, under the causal mask, lower triangular. Geometrically it is an \textbf{averaging / Markov operator on $P$}: each destination position will receive a \emph{convex combination} of source contributions. Through $S$, the operator $A$ depends on $X$ --- and this dependence (softmax of a quadratic form in $X$) is the sole source of nonlinearity in the entire head.
\end{enumerate}

\begin{enumerate}
\item \textbf{Output (move values, transform them).} With $A\in\operatorname{End}(P)$ the attention pattern, each destination collects the $A$-weighted average of values and maps it back to $V$. Coordinate-free, the head \textbf{factors across the two tensor factors} of $V\otimes P$: \[ \operatorname{head} \;=\; W_{OV} \otimes A, \qquad W_{OV} := W_O W_V : V \to V, \] with $W_{OV}$ on the \textbf{feature} factor (\emph{what} is moved and how it is transformed) and $A$ on the \textbf{position} factor (\emph{where} it moves) --- independent actions, given $A$. This is the central decomposition of \textcite{framework}. (In the $d_{\text{model}}\times n$ matrix picture, $\operatorname{head}(X)=W_{OV}\,X\,A^\top$; that transpose is \textbf{neither a dual nor an adjoint} --- only the bookkeeping of right-multiplying the position index.)
\end{enumerate}

\textbf{The two recurring circuits}, both of rank $\le d_{\text{head}}$ (they factor through $H$), and --- the point worth absorbing --- of different \emph{type}:

\begin{itemize}
\item \textbf{QK circuit} $W_{QK} : V \to V^{*}$, a \textbf{bilinear form} ($(0,2)$-tensor in $V^{*}\!\otimes V^{*}$) governing \emph{where} a head attends. It lands in the dual --- hence the transpose in the coordinate form $W_Q^{\top}W_K$ --- and feeds the nonlinear softmax.
\item \textbf{OV circuit} $W_{OV} = W_O W_V : V \to V$, an \textbf{endomorphism} (in $V^{*}\!\otimes V$) governing \emph{what} is written when a position is attended to. No transpose appears --- no dual is involved. (Its image is the head's \textbf{write subspace} of §1.3; the eigenstructure of the endomorphism $W_{OV}$, read in the token basis, diagnoses e.g. a "copying" head --- positive real eigenvalues mean attending to token $t$ raises $t$'s own logit.)
\end{itemize}

That one is a \textbf{bilinear form} and the other an \textbf{endomorphism} is the where/what distinction made structural: \emph{where} pairs two stream vectors into a number (a score), so it is a $V\to V^{*}$ object; \emph{what} sends a stream vector to a stream vector to add, so it is a $V\to V$ object. The presence of a transpose in $W_{QK}$ and its absence in $W_{OV}$ is exactly that type signature.


\paragraph*{Freezing the attention pattern --- the move that makes circuits tractable}\mbox{}\\

The head map (matrix picture) $X \mapsto W_{OV}\,X\,A(X)^\top$ is nonlinear, but in a very controlled way: \textbf{the only nonlinear dependence on $X$ is through $A = A(X)$.} Hold $A$ fixed at a constant matrix $A_0$ and the residual map \[ X \;\longmapsto\; (W_{OV}\otimes A_0)\,X \qquad(\text{matrix picture: } W_{OV}\, X\, A_0^\top) \] is \textbf{exactly linear} --- it is a single fixed linear operator on $V\otimes P$, the tensor product of two constant linear maps.


The conceptual content: the stream $X$ plays \emph{two roles} in a head. It determines the \textbf{routing} (through the QK path $X \mapsto A$) and it supplies the \textbf{content} that gets moved (through the value/OV path). "Freezing $A$" means computing the routing $A_0 = A(X_0)$ on some reference input $X_0$, then \emph{detaching} it --- regarding the head as a function of the content-carrying copy of $X$ alone, with routing held fixed. Under that split the head is the constant linear operator above.


Two properties make this powerful, and distinguish it sharply from a Taylor/Jacobian linearization:

\begin{itemize}
\item \textbf{Exact at the operating point, not first-order.} At $X = X_0$ we have $A_0 = A(X_0)$, so $W_{OV} X_0 A_0^\top$ equals the head's \emph{true} output on $X_0$ --- exactly, with no remainder. Freezing is not an approximation of the value; it is an exact rewriting that happens to be linear in the content direction. (Contrast §3's \emph{local linearization}, which matches the true map only to first order and drifts away from $X_0$.)
\item \textbf{Globally linear in content.} As a function of the value-path input (routing fixed), it is linear \emph{everywhere}, not just near $X_0$.
\end{itemize}

What you give up is precisely the dependence of routing on the input --- the derivative of $A$ through the QK circuit. That discarded piece is the genuinely nonlinear, and empirically the \emph{harder}, half of attention: explaining \textbf{why} a head attends where it does. The field handles this by \textbf{splitting the head in two} along exactly this seam --- study the OV circuit with $A$ frozen (clean linear algebra), and study the QK circuit's routing separately (still largely open; see \textcite{tracing-attention-feature-interactions}, \textcite{progress-on-attention}).


Why this "makes circuit analysis tractable": with every head's $A$ frozen to constants, an entire attention layer becomes a fixed linear operator on $V\otimes P$. Composing such layers with the (linear) residual stream of §1.3, the embedding $W_E$, and the unembedding $W_U$ yields an \textbf{end-to-end linear map from input tokens to logits}, which can be multiplied out into a sum of interpretable paths (\emph{path expansion}; \textcite{framework}, the \emph{Landscape} survey, Part I). The OV circuit then says, for the established routing, exactly which token-content is copied or transformed into which logit direction --- a question of reading eigen/singular structure of fixed matrices. The one caveat: $A$ is input-dependent, so the frozen-linear model is valid \emph{for the input it was frozen on}; on a different input the pattern, hence the linear operator, changes. (This per-input frozen-attention model is exactly the "local replacement model" of the circuit-tracing program --- §3, the \emph{Landscape} survey, Group 4.)


\paragraph*{The token-to-token circuits, and a 1-layer model attending by hand}\mbox{}\\

The path expansion of a \textbf{1-layer attention-only} model (\emph{Landscape} survey, Group~1) conjugates the two circuits above by the embedding and unembedding to get two honest $n_{\text{vocab}}\times n_{\text{vocab}}$ \textbf{token-to-token} matrices: the \textbf{QK circuit} $W_E^{\vee}\,W_{QK}\,W_E$ (coordinates $W_E^{\top}W_{QK}W_E$), scoring how much a query \emph{token} wants to attend to a key \emph{token}, and the \textbf{OV circuit} $W_U\,W_{OV}\,W_E$, recording how an attended token shifts the output logits. Each has rank $\le d_{\text{head}}$ (it factors through $H$). The worked example below instantiates both on numbers small enough to run softmax by hand; it is the concrete companion to the where/what story above.

\begin{callout}
\textbf{Worked example: a skip-trigram head.}

\emph{What the model does, and what this head is for.} The whole model is a \textbf{next-token predictor}: feed it a sequence of tokens (a \emph{prefix} of text) and it returns a probability distribution over the single token that comes next. You \emph{read} a piece of text by running the model once over it; you \emph{generate} text by sampling the predicted token, appending it to the prefix, and running again --- one token at a time, left to right. Our toy head is built to handle one scrap of English, the phrase ``\texttt{keep \dots\ in mind}''. Concretely we will hand it the prefix ``\texttt{keep this in}'' and ask it to predict the next token, which should be ``\texttt{mind}'' (the earlier tokens we simply take as given).

\emph{Three tokens, three different roles --- and ``attend'' is not ``predict''.} The step that trips up newcomers is that \textbf{three distinct tokens} are in play, doing three different jobs:
\begin{itemize}
\item the \textbf{source}, a.k.a. the \textbf{key} the head lands on: \texttt{keep}, sitting at an \emph{earlier} position --- what the head \emph{looks back at};
\item the \textbf{destination}, a.k.a. the \textbf{query}: \texttt{in}, the \emph{current} position, the one whose next token we are predicting --- what does the \emph{looking};
\item the \textbf{output}: \texttt{mind}, the token whose logit the head ultimately \emph{raises} --- the actual prediction.
\end{itemize}
So ``query \texttt{in} wants key \texttt{keep}'' means: \emph{when the current token is \texttt{in}, the head reaches backward through the context and pulls information out of the earlier position holding \texttt{keep}.} It does \textbf{not} mean the model outputs \texttt{keep}. Attention only \emph{moves information between positions}; a \emph{separate} map (the OV circuit) then turns the retrieved \texttt{keep} into the prediction \texttt{mind}. Nothing is reversed in text order: \texttt{keep} occurs first, \texttt{in} later, and \texttt{mind} is what comes after \texttt{in} --- ``\texttt{keep this in mind}''. The natural misreading --- that the head ``generates \texttt{keep} after \texttt{in}'' --- conflates the \emph{key} the head reads with the \emph{output} it writes; they are different tokens.

\emph{The setup.} Take $n_{\text{vocab}}=4$ with token basis ordered $(\texttt{keep},\texttt{in},\texttt{mind},\texttt{this})$. To strip away everything but the two circuits, identify $V\cong T$ via $W_E = I$ (each token embeds to its own basis vector) and tie the unembedding $W_U = I$ (coordinates). Then $W_UW_E=I$ --- the \emph{direct path} $\operatorname{Id}\otimes W_UW_E$ just copies the current token's one-hot to the logits (the bigram baseline) --- and the two token-to-token circuits are \emph{literally} $W_{QK}$ and $W_{OV}$ read in the token basis. Give the single head these two rank-$1$ circuits (rows/cols ordered as above):
\[
\begin{aligned}
\underbrace{W_E^{\top}W_{QK}W_E}_{\text{QK circuit}}
&=\;4\,(\mathbf e_{\texttt{in}})(\mathbf e_{\texttt{keep}})^{\top}
=\begin{pmatrix}0&0&0&0\\ 4&0&0&0\\ 0&0&0&0\\ 0&0&0&0\end{pmatrix},\\[4pt]
\underbrace{W_UW_{OV}W_E}_{\text{OV circuit}}
&=\;5\,(\mathbf e_{\texttt{mind}})(\mathbf e_{\texttt{keep}})^{\top}
=\begin{pmatrix}0&0&0&0\\ 0&0&0&0\\ 5&0&0&0\\ 0&0&0&0\end{pmatrix}.
\end{aligned}
\]
Read in those three roles: the QK row for \emph{query} \texttt{in} has its one nonzero entry in the \emph{key} column \texttt{keep} (value $4$) --- the ``\emph{when current token is \texttt{in}, look back at \texttt{keep}}'' instruction; the OV column for \emph{source} \texttt{keep} has its one nonzero entry in the \emph{output} row \texttt{mind} (value $5$) --- the ``\emph{having found \texttt{keep}, push up the prediction \texttt{mind}}'' instruction. Both are rank $1$, so $d_{\text{head}}=1$ realises them, e.g. $W_K:\texttt{keep}\mapsto 1$ (else $0$), $W_Q:\texttt{in}\mapsto 4$, and $W_V:\texttt{keep}\mapsto 1$, $W_O:1\mapsto 5\,\mathbf e_{\texttt{mind}}$.

\emph{Run it on the prefix} $\texttt{keep},\ \texttt{this},\ \texttt{in}$ (positions $1,2,3$), asking for the token after position $3$. The destination is position $3$, query token \texttt{in}; the causal mask allows keys at positions $1,2,3$, with key tokens $\texttt{keep},\texttt{this},\texttt{in}$.

\textbf{Step 1 --- scores.} Read row \texttt{in} of the QK circuit against each key token:
\[
s_1=\langle\texttt{in}\mid\texttt{keep}\rangle=4,\qquad s_2=\langle\texttt{in}\mid\texttt{this}\rangle=0,\qquad s_3=\langle\texttt{in}\mid\texttt{in}\rangle=0.
\]
\textbf{Step 2 --- pattern (the one nonlinearity).} Row-softmax over positions $1,2,3$ (folding $1/\sqrt{d_{\text{head}}}$ into the $4$):
\[
A_{3,\cdot}=\operatorname{softmax}(4,0,0)=\frac{(e^4,\,1,\,1)}{e^4+2}\approx(0.965,\;0.018,\;0.018).
\]
The destination puts $\approx\!96.5\%$ of its attention on position $1$ --- it has \textbf{attended to} \texttt{keep}.

\textbf{Step 3 --- output (OV moves the content).} The head adds $\sum_j A_{3,j}\,W_{OV}\,\mathbf e_{(\text{token }j)}$. In the token basis $W_{OV}\mathbf e_{\texttt{keep}}=5\,\mathbf e_{\texttt{mind}}$ and $W_{OV}\mathbf e_{\texttt{this}}=W_{OV}\mathbf e_{\texttt{in}}=0$, so the head's contribution to the position-$3$ logits is
\[
0.965\cdot 5\,\mathbf e_{\texttt{mind}}\;+\;0.018\cdot 0\;+\;0.018\cdot 0\;\approx\;4.82\,\mathbf e_{\texttt{mind}}.
\]
Adding the direct path ($+1$ to \texttt{in}), the logit vector is $\approx(\texttt{keep}:0,\ \texttt{in}:1,\ \texttt{mind}:4.82,\ \texttt{this}:0)$, so the model predicts \textbf{\texttt{mind}} --- completing ``\texttt{keep this in mind}''. We have read the \textbf{skip-trigram} $[\texttt{keep}]\dots[\texttt{in}]\!\to\![\texttt{mind}]$ straight off the two matrices: a nonzero QK entry (\emph{where}: \texttt{in} attends back to \texttt{keep}) composed with a nonzero OV entry (\emph{what}: an attended \texttt{keep} boosts \texttt{mind}).

\emph{Why the factoring is lossy (the "skip-trigram bug").} Suppose the same head must also do $[\texttt{keep}]\dots[\texttt{at}]\!\to\![\texttt{bay}]$. That means QK row \texttt{at} also points at \texttt{keep}, and OV column \texttt{keep} also boosts \texttt{bay}. But the head's effect on a $(\text{destination},\text{out})$ pair factors as $\big(\text{attention to source}\big)\times\big(\text{OV}[\text{out},\text{source}]\big)$ --- the destination enters only through QK, the output only through OV, coupled solely by the \emph{shared} source token. With \texttt{keep} as source attracting both destinations \emph{and} boosting both outputs, all four combinations fire, including the unwanted $[\texttt{keep}]\dots[\texttt{in}]\!\to\![\texttt{bay}]$. No choice of weights separates them: that is the documented bug, and it is visible here as ``column \texttt{keep} of the OV circuit is one vector, shared by every destination that attends to \texttt{keep}.''
\end{callout}


\textbf{Multi-head.} Several heads run in parallel per layer over the \emph{same} stream; their outputs are \textbf{summed} back in, so the layer is $\sum_h W_{OV}^h\otimes A^h$ (matrix picture $\sum_h W_{OV}^h\, X\, (A^h)^\top$). (Implementation-wise the per-head $H_h$ are concatenated and $W_O$ is correspondingly block-structured, but additivity across heads --- like additivity across blocks in §1.3 --- is the mathematically load-bearing fact.) Because each $W_{OV}^h, W_{QK}^h$ has rank $\le d_{\text{head}}$, a head reads from and writes to only a $d_{\text{head}}$-dimensional slice of the $d_{\text{model}}$-dimensional stream.


\subsubsection*{1.5 MLP (feed-forward) layer}
The other block type: a two-layer perceptron acting on the feature space $V$, \textbf{independently at each position}, through a high-dimensional \textbf{neuron space} $N = \mathbb{R}^{d_{\text{mlp}}}$ (typically $\dim N \approx 4\dim V$): \[ \mathrm{MLP} = W_{\text{out}}\circ \sigma \circ \big(W_{\text{in}}(\cdot) + b_{\text{in}}\big) + b_{\text{out}}, \qquad W_{\text{in}} : V \to N,\quad W_{\text{out}} : N \to V, \] with $\sigma$ a \textbf{pointwise} nonlinearity on $N$ (historically ReLU $\max(0,\cdot)$; modern: \textbf{GeLU} $x\,\Phi(x)$, a smooth ReLU; \textbf{SoLU} $x \odot \operatorname{softmax}(x)$, see the \emph{Landscape} survey). The coordinates of $N$ are the \textbf{neurons}. MLPs hold $\sim$2/3 of a transformer's parameters and are far less understood than attention.


Because $\sigma$ acts coordinate-wise, it is the operation that refers to --- and thereby \textbf{privileges} --- the neuron basis of $N$. Unlike the feature space $V$ (no privileged basis, §1.3), $N$ comes with a distinguished basis, which is why early interpretability hoped to read meaning off individual neurons --- and why \emph{superposition} (§3) is the obstruction when meaning fails to align with that basis.


\subsubsection*{1.6 LayerNorm / RMSNorm}
A near-normalization of the stream vector, applied per position before a block reads it; this is the one place the coordinates of $V$ are genuinely referenced.

\begin{itemize}
\item \textbf{LayerNorm}: $\operatorname{LN}(x) = \gamma \odot \dfrac{x - \mu(x)\,\mathbf 1}{\varsigma(x)} + \beta$, where $\mu(x)$ and $\varsigma(x)$ are the mean and standard deviation of the coordinates of $x$, and $\gamma,\beta$ are learned. Two steps are basis-dependent: mean-subtraction is the orthogonal projection \emph{off the all-ones line} $\mathbb{R}\mathbf 1$ (privileging that direction), and the learned $\gamma\odot(\cdot)$ is a fixed diagonal map in the coordinate basis. The sole nonlinearity is the division by the scalar $\varsigma(x) = \|x - \mu(x)\mathbf 1\|/\sqrt{\dim V}$, which is often "folded in" or \emph{frozen} (treated as a constant) during analysis, leaving LN affine.
\item \textbf{RMSNorm}: drops the mean-centering --- $x \mapsto \gamma \odot x / \operatorname{RMS}(x)$, $\operatorname{RMS}(x)=\|x\|/\sqrt{\dim V}$. Apart from the learned diagonal $\gamma$, this depends on $x$ only through $\|x\|$ and the direction $x/\|x\|$, hence is \textbf{rotation-invariant} --- whereas LayerNorm's mean-subtraction privileges $\mathbf 1$. Used in experiments that need exact basis-independence (e.g. the rotation tests in \textcite{privileged-bases-residual-stream}).
\end{itemize}

\subsubsection*{1.7 Positional information}
Tokens have no intrinsic order, so position is injected: learned \textbf{positional embeddings} added to the stream, or \textbf{RoPE} (rotary positional embeddings), which rotates $q,k$ by a position-dependent angle --- meaning some positional information lives \emph{in the QK bilinear form}, not in the residual stream (a caveat for attribution methods).



\subsection*{2. Training}

\begin{itemize}
\item \textbf{Loss}: training minimizes \textbf{cross-entropy} $-\log \mathbb{P}_{\text{model}}(\text{actual next token})$, averaged over a corpus. Reported in \textbf{nats} (natural-log units). Equivalent to maximizing log-likelihood; the model is a density estimator over text.
\item \textbf{Gradient descent / SGD}: parameters updated by $-\eta \nabla_\theta L$ on mini-batches.
\item \textbf{Adam / AdamW}: the optimizer of choice. It keeps \textbf{per-coordinate} running estimates of the gradient's first and second moments and normalizes the update coordinate-wise. The "\textbf{W}" adds decoupled \textbf{weight decay} (an $\ell_2$ shrinkage of weights each step). \emph{Per-coordinate normalization is basis-dependent} --- this is the suspected cause of basis privilege in the residual stream (see the \emph{Landscape} survey).
\item \textbf{Weight decay / regularization}: penalties added to the loss to bias toward small/simple weights --- $\ell_2$ (ridge), $\ell_1$ (sparsity-inducing), or $\ell_0$ (literal count of nonzeros, NP-hard, approximated by $\ell_1$). Central to dictionary learning (SAEs) and to grokking.
\item \textbf{Epoch / full-batch / learning-rate schedule}: an epoch is one pass over the data; toy models often use full-batch (exact gradient); learning rates are typically warmed up then decayed.
\item \textbf{Pretraining vs finetuning}: \emph{pretraining} is the big next-token run on web-scale text; \emph{finetuning} (incl. \textbf{SFT}, \textbf{RLHF} --- reinforcement learning from human feedback) adapts the pretrained model to be a helpful assistant, refuse harmful requests, etc.
\item \textbf{Scaling laws}: empirically, loss falls as a power law in compute/parameters/data. Used here to budget SAE training too.
\end{itemize}

\subsubsection*{2.1 Two training phenomena that are themselves objects of study}
\begin{itemize}
\item \textbf{Grokking}: on small algorithmic tasks, a network first \textbf{memorizes} the training set (train loss $\to 0$, test loss high), then --- often long after --- \textbf{suddenly generalizes} (test loss drops). See the \emph{Landscape} survey, Part II, Group 5, and \textcite{progress-measures-for-grokking}.
\item \textbf{Double descent}: test error, as a function of model size or dataset size, can go down, then \textbf{up} (around the interpolation threshold), then down again --- contradicting the classical U-shaped bias--variance picture. See \textcite{superposition-memorization-double-descent}.
\end{itemize}


\subsection*{3. Interpretability-specific vocabulary}

This is the vocabulary invented by the literature itself; the survey's mathematical arc is mostly about these.


\begin{itemize}
\item \textbf{Feature}: the unit of "meaning." The dominant working definition (the \textbf{Linear Representation Hypothesis}) is that a feature is a \textbf{direction} (a unit vector, or a ray/cone) in some activation space, and that the activation decomposes as a sparse linear image $x \approx b + D\,f(x)$ of a nonnegative sparse code $f(x)\in\mathbb{R}^F_{\ge 0}$ under a \textbf{dictionary} map $D : \mathbb{R}^F \to V$ whose values $d_i := D e_i$ on the standard basis are the feature directions (equivalently $x \approx b + \sum_i f_i(x)\, d_i$). A more recent minority view (\texttt{when-models-manipulate-manifolds}) is that a feature can be a \textbf{curved low-dimensional manifold}, not a single direction. See the \emph{Landscape} survey, Part IV.
\item \textbf{Activation}: the actual numerical vector at some site (residual stream at layer $k$, MLP neurons, attention values) for a given input. Contrast \textbf{weights} (fixed parameters).
\item \textbf{Monosemantic vs polysemantic}: a neuron/feature is \emph{monosemantic} if it responds to one human-interpretable concept, \emph{polysemantic} if it fires for several unrelated ones. Polysemanticity is the empirical norm for neurons and is the symptom of superposition.
\item \textbf{Privileged basis}: a basis of an activation space that is \emph{singled out by the architecture}. Pointwise nonlinearities (ReLU on MLP neurons) create one; purely linear surroundings (the residual stream, attention value space) do not. Formally: a representation space $U$ is \emph{non-privileged} if inserting $M$ before and $M^{-1}$ after leaves the function unchanged for any automorphism $M\in GL(U)$ --- i.e. it has a \textbf{$GL(U)$ gauge symmetry}. This gauge idea is a unifying thread in the \emph{Landscape} survey.
\item \textbf{Superposition}: a layer of width $d$ representing $\gg d$ features as an \textbf{overcomplete} set of almost-orthogonal directions, made recoverable by \textbf{sparsity} (few features active at once). The enabling mathematics: \textbf{Johnson--Lindenstrauss} (see §4) and \textbf{compressed sensing}. For a linear code given by an embedding map $W : e_i \mapsto d_i$ (feature-coordinate space $\to$ the width-$d$ space), superposition is present \textbf{iff the feature directions $\{d_i\}$ are linearly dependent} --- equivalently $W$ is non-injective ($\ker W\neq 0$), a metric-free condition; equivalently (invoking the metric) $W^{*}W$ is singular --- which is forced once there are more features than dimensions. See \textcite{toy-models-superposition}.
\item \textbf{Interference}: the off-target inner products $\langle d_i, d_j\rangle$, $i\ne j$, between feature directions --- the "noise" superposition pays for packing. The off-diagonal of the Gram matrix $D^{*}D$ (the adjoint $D^{*}$ makes the metric-dependence explicit; §4). "\textbf{Interference weights}" are the spurious effective weights this induces between features (\textcite{toy-model-of-interference-weights}).
\item \textbf{Sparse Autoencoder (SAE)}: the workhorse tool. A wide autoencoder reconstructing activations $x\in V$ through an \textbf{overcomplete} ($F \gg \dim V$) sparse code in $\mathbb{R}^F$:
\begin{itemize}
\item encoder $f = \sigma\big(W_{\text{enc}}(x) + b_{\text{enc}}\big)$, with $W_{\text{enc}} : V \to \mathbb{R}^F$ ($\sigma=$ ReLU or \textbf{JumpReLU}),
\item decoder $\hat x = b_{\text{dec}} + D\,f(x)$, with the \textbf{dictionary} map $D = W_{\text{dec}} : \mathbb{R}^F \to V$,
\item loss $\mathbb{E}\big[\;\|x-\hat x\|_2^2 + \lambda \sum_i f_i(x)\,\|D e_i\|_2 \big]$ (reconstruction + decoder-norm-weighted $\ell_1$ sparsity). The images $D e_i$ are the recovered \textbf{feature directions}; this is \textbf{dictionary learning} with a learned overcomplete dictionary $D$. See the \emph{Landscape} survey, Part II, Group 3.
\end{itemize}
\item \textbf{Dictionary learning / sparse coding}: classical problem (Olshausen--Field, k-SVD): given data, find an overcomplete dictionary $D$ and sparse codes such that $x \approx D f$, $f$ sparse. SAEs are a neural, amortized version.
\item \textbf{Transcoder}: like an SAE, but instead of reconstructing its own input it predicts a \emph{later} activation (e.g. an MLP's \textbf{output} from its \textbf{input}) --- so it models \textbf{computation}, not just representation. \textbf{Cross-layer transcoder (CLT)}: features at layer $\ell$ write to all later layers. \textbf{Crosscoder}: one shared dictionary reading/writing multiple sites (layers, or even different models) at once.
\item \textbf{JumpReLU}: a thresholded ReLU ($0$ below a learned threshold, identity above) used to get crisper sparsity.
\item \textbf{Dead feature}: an SAE feature that (almost) never activates --- wasted capacity; mitigated by "resampling."
\item \textbf{Feature splitting}: as dictionary width $F$ grows, one coarse feature resolves into several finer ones (e.g. "San Francisco" $\to$ many SF-related features). Suggests a hierarchy / no canonical feature count.
\item \textbf{Probe (linear probe)}: a linear (or logistic) classifier given by a \textbf{covector} $w\in V^{*}$, reading off the canonical pairing $\langle w, x\rangle = w(x)$ --- no metric needed. Tests whether information is \textbf{linearly decodable}. A "concept direction" is properly this covector $w$; calling it a \emph{direction} (a vector in $V$) silently transports it across the metric isomorphism $V^{*}\cong V$.
\item \textbf{Steering / activation steering / clamping}: \emph{intervening} at inference by adding a vector ($x \mapsto x + \alpha v$) or fixing a feature's value, to causally change behavior. Tests whether a direction is causal, not merely correlational.
\item \textbf{Circuit}: a subgraph of the model's computation (specific heads/features and the connections among them) implementing a specific behavior. The "circuits" program reverse-engineers these.
\item \textbf{Ablation}: deleting a component (set its output to \textbf{zero} --- \emph{zero ablation}, or to its dataset \textbf{mean} --- \emph{mean ablation}) and measuring the behavioral change. A causal-importance test.
\item \textbf{Activation patching / path patching / attribution patching}: causal-intervention techniques. \emph{Activation patching} copies an activation from a "clean" run into a "corrupted" run (or vice versa) and measures the effect. \emph{Path patching} restricts this to specific paths between components. \emph{Attribution patching} is the cheap \textbf{linear (gradient) approximation} of patching: effect $\approx \nabla_x(\text{metric})\cdot \Delta x$.
\item \textbf{Attribution graph}: a per-prompt causal graph (nodes = features/tokens/logits, edges = linear attributions) extracted from a locally-linearized replacement model. See \textcite{circuit-tracing-computational-graphs,on-the-biology-of-a-large-language-model}.
\item \textbf{Faithfulness / completeness / minimality}: the validation criteria for a claimed circuit $C$ relative to the full model $M$ (formalized in \textcite{interpretability-in-the-wild-ioi}). Roughly: \emph{faithful} = $C$ alone reproduces $M$'s behavior; \emph{complete} = $C$ contains all the relevant components; \emph{minimal} = no component of $C$ is redundant. \textbf{Mechanistic faithfulness} (a distinct, deeper notion, \textcite{toy-model-of-mechanistic-unfaithfulness}) = the approximation implements the \emph{same algorithm}, not merely the same input--output function.
\item \textbf{In-context learning (ICL)}: the ability to learn a pattern from the prompt itself (no weight updates), e.g. few-shot examples. Largely attributed to \textbf{induction heads}.
\item \textbf{Induction head}: a two-head circuit implementing the copy-completion rule $[A][B]\dots[A]\to[B]$ (and fuzzy versions). The canonical worked example of the whole framework. See \textcite{induction-heads}.
\end{itemize}


\subsection*{4. Mathematical facts the papers lean on (you know these, but here's the ML framing)}

\begin{itemize}
\item \textbf{Johnson--Lindenstrauss lemma}: in $\mathbb{R}^{d}$ one can place \textbf{exponentially many --- $e^{\Omega(\varepsilon^2 d)}$ --- unit vectors that are pairwise almost orthogonal} ($|\langle u_i,u_j\rangle| < \varepsilon$). (Equivalently: any finite set of points embeds linearly into dimension $O(\varepsilon^{-2}\log(\#\text{points}))$ with pairwise distances preserved to $1\pm\varepsilon$.) This is the existence proof for superposition --- the feature space $V$ has room for vastly more than $\dim V$ nearly-orthogonal feature directions.
\item \textbf{Compressed sensing}: a $k$-sparse vector in $\mathbb{R}^{F}$ is recoverable from $\Omega\!\big(k\log(F/k)\big)$ linear measurements, provided the measurement map is "incoherent" (RIP / restricted isometry; low \textbf{coherence} $\max_{i\ne j}|\langle d_i,d_j\rangle|$). Read into the architecture: a model can pack $F$ sparse features (at most $k$ active at once) into a feature space of dimension $\dim V \gtrsim k\log(F/k)$ and still recover them --- and an SAE is the recovery algorithm. This \emph{upper-bounds} how much superposition is possible.
\item \textbf{Coherence / Gram matrix}: for a dictionary map $D:\mathbb{R}^F\to V$ with unit-norm columns $d_i = D e_i$, the \textbf{Gram matrix} $G = D^{*}D$ (an endomorphism of the code space $\mathbb{R}^F$, entries $G_{ij} = \langle d_i, d_j\rangle$) is built from the \textbf{adjoint} $D^{*}$, so it depends on the chosen inner product on $V$. Its off-diagonal is interference, its diagonal is $1$; near-orthogonality means $G \approx I$. Almost every quantitative claim in this literature is, under the hood, a statement about $G$ --- and hence about that silent choice of metric (see the \emph{Landscape} survey, Part III).
\item \textbf{Thomson problem / packing}: minimizing pairwise interaction energy of points on a sphere; its discrete optimal configurations are \textbf{uniform polytopes} (simplices, antiprisms, \dots{}). \textcite{toy-models-superposition} shows learned feature geometries \emph{are} such configurations.
\item \textbf{Discrete Fourier transform / characters of $\mathbb{Z}/n\mathbb{Z}$}: the irreducible representations of the cyclic group are $k \mapsto e^{2\pi i k j/n}$. A network doing modular addition represents inputs on these circles and adds angles (\textcite{progress-measures-for-grokking}). The same circular/periodic ("ringing") structure resurfaces geometrically in \texttt{when-models-manipulate-manifolds}.
\item \textbf{Riesz representation}: every continuous linear functional on a Hilbert space is $\langle v,\cdot\rangle$ for a unique $v$; the map $v\mapsto\langle v,\cdot\rangle$ is the \textbf{musical isomorphism} $V\xrightarrow{\sim}V^{*}$ induced by the inner product --- the same identification that turns a \textbf{dual map} $W^{\vee}$ into an \textbf{adjoint} $W^{*}$. \texttt{linear-representation-hypothesis} uses exactly this to identify the embedding-space and unembedding-space representations of a concept under a chosen (non-Euclidean) inner product.
\end{itemize}


\emph{Append new concepts below as later surveys introduce them.}



\printbibliography
\end{document}
